\section{Worked Example: A Specialized Technical Training Market}

To see how bottom-up sizing works in practice, examine a concrete hypothetical scenario for a premium artificial intelligence engineering program in a major Indian technology hub.

Assume the region contains \textbf{60,000 employed software engineers} across all experience levels. To construct a disciplined model, apply explicit, falsifiable behavioral filters rather than wishful thinking:

\begin{tcolorbox}[
  enhanced,
  colback=BayesSurface!60,
  colframe=BayesBorder,
  boxrule=0.5pt,
  arc=3pt,
  left=3pt,right=3pt,top=3pt,bottom=3pt,
  boxsep=1pt
]
\scriptsize
\begin{tabularx}{\linewidth}{>{\bfseries\raggedright\arraybackslash}p{2.3cm} >{\raggedleft\arraybackslash}p{1.4cm} X}
\toprule
\textbf{Analytical Filter} & \textbf{Modeled Count} & \textbf{Operational Interpretation} \\
\midrule
Total Tech Professional Population & 60,000 & Background technical labor pool; \textbf{not} commercial demand. \\
\addlinespace[2pt]
Prerequisite Eligibility (20\%) & 12,000 & Developers possessing required backend math and coding fundamentals. \\
\addlinespace[2pt]
Annual Triggering Need (25\%) & 3,000 & Engineers tasked with an immediate, acute production AI deployment. \\
\addlinespace[2pt]
Willingness and Ability to Pay (30\%) & 900 & Decision-makers capable of committing personal or employer funds. \\
\addlinespace[2pt]
Standard Program Tuition & ₹40,000 & Realistic realized fee for an intensive 8-week applied clinic. \\
\midrule
\textbf{Annual Serviceable Market (SAM)} & \textbf{₹3.6 crore} & $900 \times \text{₹40,000}$; total theoretical category spend, \textbf{not} startup revenue. \\
\bottomrule
\end{tabularx}
\end{tcolorbox}

The resulting ₹3.6 crore figure represents the total addressable pool of capital circulating across all competitors, self-study resources, and employer programs for this specific niche. It does not mean our startup will capture ₹3.6 crore.

Now, construct the startup's real \textbf{Serviceable Obtainable Market (SOM)} based on verified marketing channels and operational capacity:
\begin{enumerate}[leftmargin=12pt,itemsep=1pt,topsep=1pt,parsep=0pt]
  \item Through direct outbound technical publishing and developer communities, the founder can reach \textbf{180 qualified prospects} in Year 1.
  \item Diagnostic qualification filters narrow this group to \textbf{60 serious applicants}.
  \item Substantive technical evaluation interviews produce \textbf{30 qualified offers}.
  \item A realistic 40\% paid conversion yield produces \textbf{12 enrolled learners}.
\end{enumerate}

At ₹40,000 per seat, this yields \textbf{₹4.8 lakh in gross tuition revenue in Year 1.}

Notice the crucial insight: if the founder can responsibly mentor only 8 learners simultaneously without compromising instructional quality, \textbf{the first-year business is bounded by delivery capacity (8 seats = ₹3.2 lakh), not by market demand.}

\section{Worked Example: A College Partnership Market}

The market for an institutionally licensed college AI elective cannot be estimated by multiplying India's 4.5 crore university students by a course fee. The college administration is the buyer, students are users, faculty committees enforce academic compliance, and contracts are negotiated per department or academic year.

Start with a \textbf{named, verified target list} of institutions within accessible geographic travel range:
\begin{enumerate}[leftmargin=12pt,itemsep=1pt,topsep=1pt,parsep=0pt]
  \item The founder identifies \textbf{40 accredited private engineering colleges} seeking to modernize their data science curricula.
  \item Exploratory meetings with academic deans yield \textbf{10 institutions} interested in an external curriculum partnership.
  \item Formal curricular reviews and faculty alignment narrow the pipeline to \textbf{4 colleges} requesting a priced pilot proposal.
  \item Exactly \textbf{1 institution} successfully navigates governing board approval, signs an academic service agreement, and issues a purchase order.
\end{enumerate}

At an agreed pilot fee of ₹3,00,000 for a single semester cohort of 60 students, the startup books \textbf{₹3,00,000 in contracted gross revenue}.

However, until the course is delivered, proctored assessments are completed, and accounts payable releases funds (often 60 to 90 days after delivery), \textbf{contracted bookings are not collected cash.} Fulfillment costs, travel expenses, teaching assistant stipends, and collection delays determine whether the engagement generates positive contribution margin.

\section{Worked Example: An Enterprise AI Capability Service}

Now consider an enterprise business model: providing a specialized evaluation-and-training engagement to engineering teams deploying generative AI pipelines.

Here, the buyer unit is an \textbf{engineering team with independent budget authority and a production AI deployment roadmap}, not an individual corporate employee.

A credible bottom-up calculation proceeds as follows:
\begin{itemize}[leftmargin=12pt,itemsep=1pt,topsep=1pt]
  \item Identify 150 target enterprises operating in fintech, healthcare, and logistics within regional technology corridors.
  \item Filter by enterprises with active customer-facing AI applications (estimated at 30\% = 45 companies).
  \item Account for internal procurement structures: some enterprises centralize training under HR learning budgets, while others allow engineering directors to expense technical consulting.
  \item The founder can secure exploratory discovery meetings with 12 engineering directors.
  \item Detailed diagnostic scoping leads to 3 formal enterprise pilot proposals.
  \item Exactly 1 enterprise approves a ₹5,00,000, 4-week intensive diagnostic engagement.
\end{itemize}

The critical operational takeaway: \textbf{a small founding team can sell and fulfill only 1 or 2 high-touch enterprise engagements per quarter.} Even if a theoretical enterprise market contains thousands of teams, the founder's sales pipeline and delivery bandwidth establish a hard ceiling on Year 1 revenue.
